
\section{Related Work}

\subsection{Multi-Label Image Recognition}

Unlike single-label recognition, multi-label image recognition allows one image to be assigned to multiple non-mutually exclusive categories, and therefore requires careful handling of label co-occurrence, label relevance, and class imbalance \cite{_Ref213949872}. Early methods such as Binary Relevance \cite{_Ref213949896} usually decompose a multi-label task into multiple independent binary classification problems. Such methods tend to ignore label relevance and may lead to incomplete or inconsistent predictions in complex visual scenarios.

To alleviate these problems, recent multi-label image recognition methods have increasingly focused on label relevance modeling and sample imbalance. ML-GCN \cite{_Ref213949382} constructs a label graph from label co-occurrence statistics and propagates label relevance information through graph convolutional networks. Asymmetric Loss (ASL) \cite{_Ref213949927} introduces different focusing parameters for positive and negative samples and suppresses the loss contribution of easy negative samples through negative probability shifting, thereby alleviating positive-negative sample imbalance in multi-label recognition. Query2Label \cite{_Ref213950246} further uses a Transformer query mechanism to model label interactions and achieves strong performance on natural image datasets such as VOC \cite{_Ref213950282}. These studies show that explicitly or implicitly incorporating label relevance into the model can improve the completeness and accuracy of multi-label prediction.

In medical fundus images, multi-label recognition is clinically important because multiple diseases may occur simultaneously in the same eye, and ignoring disease co-occurrence relationships may cause false positives and false negatives \cite{_Ref213949362}. Existing studies have explored multi-label ocular disease recognition from different perspectives. For example, Cheng et al. \cite{_Ref213950349} transferred ML-GCN to multi-label recognition of fundus lesions and modeled lesion label relevance by constructing a fundus lesion corpus and GloVe label representations. DCNet, proposed by He et al. \cite{_Ref213950338}, captures pixel-level relevance between left-eye and right-eye fundus image features through a spatial correlation module and fuses binocular features for patient-level multi-label ocular disease recognition. DKCNet \cite{_Ref213950343} uses an attention module to generate discriminative regional features and combines channel recalibration with resampling strategies to alleviate class imbalance in ocular datasets.

Overall, existing methods have partially addressed label relevance modeling and class imbalance. However, in fundus images with multi-disease co-occurrence, disease relationships are complex and lesion manifestations are often localized. Relying only on disease co-occurrence statistics or generic attention mechanisms may still be insufficient to characterize label-feature relevance reliably. \revision{Even when a model learns these associations in the source domain, differences in imaging conditions may alter the visibility of local disease signs, making the visual cues associated with a given label less reliable at another institution. Therefore, an important research question is how a model trained on a single source domain can learn label-feature relationships that remain useful for recognizing each disease across imaging conditions.} Building on label relevance modeling, this study couples adaptive label-feature relevance construction with single-source domain-invariant feature learning to address disease co-occurrence and cross-dataset application needs.

\subsection{Domain Generalization}

In computer vision tasks, \revision{a model may perform well on training data, but its recognition performance may decline on data acquired at other institutions or with different devices} because of differences in imaging conditions, illumination environments, and data distributions. This domain shift is especially prominent in medical imaging, where differences among hospitals, devices, and acquisition procedures can directly affect model generalization \cite{_Ref213949439}. \revision{The survey by Alzubaidi et al. \mbox{\cite{_RefAlzubaidi2023DataScarcity}} also examines generalization challenges under data scarcity and reviews learning strategies such as transfer learning and self-supervised learning.} Domain Adaptation (DA) is a common approach for alleviating domain shift, with the core idea of aligning distributions using both source-domain and target-domain data during training. For example, maximum-mean-discrepancy-based Deep Domain Confusion \cite{_Ref213949448} and the correlation alignment method CORAL \cite{_Ref213949453} improve transfer performance by reducing feature distribution differences. DANN \cite{_Ref213949479} uses adversarial training to learn domain-invariant features. CycleGAN \cite{_Ref213949490} and related unsupervised domain adaptation methods \cite{_Ref213949530} alleviate domain shift through image style transfer or cross-domain contrastive learning. However, DA usually requires access to target-domain data during training, an assumption that often fails when medical data are privacy-sensitive or data from the target institution cannot be obtained in advance.

\revision{Unlike DA, Domain Generalization (DG) aims to learn predictive models using only source-domain data for application to unseen target domains} \cite{_Ref213949554}. \revision{Domain-invariant representation learning is one approach to this task, aiming to reduce the influence of domain-related variations on the learned features.} Typical DG methods include MixStyle \cite{_Ref213949562}, which simulates potential domain shift by mixing instance style statistics; DomainBed \cite{_Ref213949607}, which provides a unified benchmark for evaluating DG algorithms; FACT \cite{_Ref213949788}, \revision{which combines Fourier-based augmentation with prediction consistency between original and augmented images to encourage stable semantic representations}; RSC \cite{_Ref213949796}, which suppresses dominant features to force the model to attend to more class-relevant cues; and meta-learning-based DG methods \cite{_Ref213949803}, which improve generalization by simulating differences between source and target domains. Most DG studies use multiple labeled source domains to learn cross-domain stable representations. This helps models capture richer inter-domain variations, but also requires consistent label spaces, coordinated data processing, and cross-institutional collaboration. In real-world medical scenarios, however, these conditions are not always available, making it difficult to centralize or synchronously use multi-institutional data. Even when centralized storage is avoided through federated learning, additional complexity arises from communication, privacy mechanisms, \revision{data that are not independent and identically distributed}, and collaborative optimization.

\revision{Against this background, Single-Domain Generalization (SDG) is a more constrained generalization setting in which a model is trained using data from a single source domain for application to unseen target domains.} Studies on single-source DG in medical images \cite{_Ref213949863} also indicate that mining more stable disease-related features from a single source domain has practical value when multi-institutional data cannot be centrally used. In summary, the assumptions of target-domain accessibility or multi-source collaboration presumed by common cross-domain methods are difficult to satisfy simultaneously in real-world medical scenarios. \revision{Therefore, learning transferable disease-related representations from a single source domain} is well aligned with the deployment constraints of cross-institutional multi-label ocular disease recognition.

\revision{In this setting, images from the target institution are unavailable during training, so the model cannot directly observe how imaging conditions at that institution affect the appearance of disease signs. Augmenting source images offers a way to simulate appearance changes during training. For example, FACT \mbox{\cite{_Ref213949788}} uses Fourier-based augmentation and encourages consistent predictions for the original and augmented images. In multi-label ocular disease recognition, however, similar scores for a given disease in the two views may be based on different visual cues. An important question, therefore, is how to learn features that remain stable under appearance changes while retaining the information needed to recognize each coexisting disease. To this end, our method filters features across the two views, uses disease-label queries to extract discriminative representations from the retained features, and encourages the attention vectors for the same disease label to remain consistent across views.}

\subsection{Privacy Preservation}

Medical images and electronic health records contain highly sensitive personal information, and model training, parameter sharing, or cloud inference may all introduce privacy leakage risk \cite{_Ref213950358}. In medical artificial intelligence applications, privacy preservation is tied not only to patient rights but also to data compliance requirements. Common privacy-preserving approaches include federated learning, encrypted computation, secure multi-party computation, and differential privacy. Among them, Federated Learning (FL) allows multiple clients to train models locally and upload parameters or gradients, thereby reducing the need for centralized sharing of raw data \cite{_Ref213950376}. For example, FedProx \cite{_Ref213950391} handles heterogeneous client data by adding a proximal term, alleviating instability in collaborative optimization.

Beyond federated learning, \revision{encrypted computation can protect data during processing}. Homomorphic Encryption (HE) allows computation over ciphertexts \cite{_Ref213950427} and has been explored for medical image analysis \cite{_Ref213950433}, but it usually incurs high computational overhead. Secure Multi-Party Computation (SMPC) \cite{_Ref213950438} can protect the input privacy of each party during multi-party collaboration. Differential Privacy (DP) injects noise into query results or training gradients to provide quantifiable constraints on the privacy leakage risk of individual samples \cite{_Ref213950443}. Among DP mechanisms, Gaussian differential privacy \cite{_Ref213950449} injects Gaussian noise into clipped gradients and provides quantifiable privacy preservation for continuous gradient updates, making it well suited to gradient protection in deep learning.

Existing privacy-preserving mechanisms can reduce privacy leakage risk, but they still face additional challenges in domain generalization tasks. Federated learning, homomorphic encryption, and secure multi-party computation often introduce communication, computation, or system-level collaboration costs. Differential privacy is relatively straightforward to deploy, but noise injection may affect model accuracy, and this effect may vary across classes or data distributions \cite{_Ref213950458}. These issues make privacy preservation more challenging under multi-domain generalization. In summary, the privacy sensitivity of medical data further restricts cross-institutional data sharing and \revision{makes the conditions for multi-source collaboration harder to meet}. \revision{For cross-institutional multi-label ocular disease recognition, the challenge is therefore to limit patient data exposure by reducing the need for data access and interaction across institutions during cross-domain learning, while protecting information about individual training samples and preserving recognition performance in unseen domains as far as possible. To address this challenge, we first adopt single-source domain generalization and restrict training updates to the available source domain. This data isolation limits exposure by reducing the need to obtain target-institution data for model adaptation or to pool raw images across institutions. Gradients and model parameters from local training may still carry information about source-domain samples. We therefore incorporate AdaGDP into the joint training of DFE and ARC to protect source-domain training information through gradient clipping and scheduled noise injection. We assess how much recognition utility this configuration retains by examining performance alongside the recorded privacy budget.}
